\chapter{All even interior exponents: complete source templates and
budget-preserving root reduction}
{\small\sffamily Source: \nolinkurl{repository/collatz_reconstruction/research_program/all_even_join_extension_20260916/note.md}.\par}


16 September 2026. Continuation from the original bilateral-root
calculation, whose source has now been integrated into \nolinkurl{main}
at \nolinkurl{9453b5306b3a897389dece309f9c119dffbcc489} in
\nolinkurl{KokunoYumeto/collatz-workbench}. The earlier restriction to
interior exponents 2 and 6 is removed on both displayed template types.
All positive even exponents are admitted. The induced ternary order,
exact integer source, and two original paths remain specified.

This note proves a larger family of comparisons, not global Collatz
convergence. A shared-future comparison identifies two original
first-jet classes; it is not necessarily forward iteration from the
larger source to the smaller source. The complete even-exponent test is
finite on every original input. One version preserves the preceding
subquadratic support budget by testing actual path support; another
admits the entire enlarged template with its own proved budget. The old
support bound is not silently assigned to the unrestricted version.

\section{1. Original source and first-jet receiver}

Let \(X=\{1,3,5,\ldots\}\), and use
\SourceMath{T(n)=\frac{3n+1}{2^{\nu_2(3n+1)}}.}{} An exponent word
\(w=(a_1,\ldots,a_j)\) retains all its ordered exponents, prefix sums
\(A_i\), and original affine data
\SourceMath{L_w=3^j,\qquad U_w=2^{A_j},\qquad
C_w=\sum_{i=0}^{j-1}3^{j-1-i}2^{A_i},\qquad
f_w(n)=\frac{L_wn+C_w}{U_w}.}{} Its complete positive source is the
odd-terminal congruence
\SourceMath{L_w\rho_w+C_w\equiv U_w\pmod{2U_w},\qquad 0<\rho_w<2U_w,}{}
with progression and image
\SourceMath{\rho_w+2U_wt\longmapsto\eta_w+2L_wt,\qquad
\eta_w=(L_w\rho_w+C_w)/U_w,\quad t\ge0.}{} The original predecessor
supplies the exact word/cylinder proof. All valuations in the new
families are also derived independently below.

Keep the absolute original graph, including the fixed vertex and loop at
1. Over \(\mathbb D=\mathbb Z[\epsilon]/(\epsilon^2)\), put
\(q=1+\epsilon\) and
\SourceMath{d_\epsilon E_n=V_n-qV_{T(n)},\qquad q^k=1+k\epsilon\quad(k\in\mathbb Z).}{}
The last identity follows for nonnegative k by the binomial formula and
for negative k from the integral inverse \(q^{-1}=1-\epsilon\). No
integer period is inverted. For actual finite paths \(l:n\to y\) and
\(r:m\to y\), of lengths s,t, retain the weighted edge columns
B\_l,B\_r. Their literal boundary identity is
\SourceMath{d_\epsilon\bigl(B_l(n)-q^{s-t}B_r(m)\bigr)=V_n-q^{s-t}V_m.}{1}
Multiplying by epsilon identifies the original reduction-kernel classes
\(\xi_n=[\epsilon V_n]\) and \(\xi_m\). It does not certify their
vanishing. Every support label below is the set of original equations
used, with its full coefficient module and incoming relation map. It is
not an automata word label.

\section{2. The previously omitted ternary orders}

For every even integer e\textgreater=2 define
\SourceMath{h_e=\frac{2^e+2}{3},\qquad t_e=\nu_3(h_e).}{} The numerator
is divisible by 3, h\_e is positive, and h\_e=2 modulo 4.

\textbf{Lemma 1.} For every such e,
\SourceMath{\boxed{t_e=\nu_3(e-1).}}{2}

\textbf{Proof.} Set k=e-1, which is odd. Then \(2^e+2=2(1-(-2)^k)\). For
k coprime to 3, expansion of \((1-3)^k\) modulo 9 gives \(1-(-2)^k=3k\)
modulo 9, with valuation one. For any integer x=1 modulo 3, write
x=1+3z. Then \(x^2+x+1=3(1+3z+3z^2)\) has valuation one. Factoring
\(x^3-1=(x-1)(x^2+x+1)\) therefore increases the valuation by exactly
one on cubing. Write k=3\^{}r k\_0 with 3 not dividing k\_0 and iterate
the factorization. It gives \(\nu_3(1-(-2)^k)=1+r\). Division by 3 in
h\_e proves (2). QED.

In particular e=4 contributes t\_e=1 and e=10 contributes t\_e=2.
Treating their h\_e as ternary units would put the family on the wrong
original source layer. The earlier values e=2,6 both have t\_e=0 and are
recovered by (2).

\section{3. One exact family for every even exponent and tail length}

Take b\textgreater=1, even e\textgreater=2, and a selector
\(\sigma\in\{0,1\}\). Define
\SourceMath{a=a(b,e,\sigma)=\min\{a\ge1:2^{a+e+2b-\sigma}<3^{a+b}\},
\quad J=2^{e+2b-\sigma},\quad Q=3^{a+b},\quad K=2^aJ<Q.}{3} The minimum
exists. At \(a=b+2e-2\sigma\), the ratio K/Q equals
\((8/9)^{b+e-\sigma}<1\), so \SourceMath{a\le b+2e-2\sigma.}{4} Also
a\textgreater=e. Indeed a\textless=e-1 would give \((3/2)^a<2^{e-1}\),
whereas (3) requires \((3/2)^a>2^{e-\sigma}(4/3)^b>2^{e-1}\). Since
h\_e\textless3\^{}e, it follows that a\textgreater t\_e.

For sigma=0 retain left word (1) and the residue n=3 modulo 4. For
sigma=1 retain left word (1,2,1) and the residue n=27 modulo 32. Write
these dyadic data as \((d_\sigma,r_\sigma)=(4,3)\) or \((32,27)\). The
complete larger-source domain is
\SourceMath{\mathcal E_{b,e,\sigma}=\{n>0:n\equiv r_\sigma\pmod{d_\sigma},
\quad Jn+h_e3^b\equiv0\pmod Q\}.}{5} Because J is a unit modulo Q, the
second condition is one residue modulo Q. CRT with the dyadic condition
has one representative n\_0 in (0,d\_sigma Q).

Put \SourceMath{m_0=\frac{Kn_0+2^ah_e3^b}{Q}-1,
\qquad y_0=f_{l_\sigma}(n_0).}{} Then the complete pair and its common
image are \SourceMath{\boxed{\begin{aligned}
n(v)&=n_0+d_\sigma Qv,\\
m(v)&=m_0+d_\sigma Kv,\\
y(v)&=y_0+\frac{L_{l_\sigma}d_\sigma Q}{U_{l_\sigma}}v,
\qquad v\in\mathbb Z_{\ge0}.
\end{aligned}}}{6} The original right word is
\SourceMath{w_0=1^a,e,2^{b-1},3,\qquad
w_1=1^a,e,2^{b-1},1,1,3.}{7} Here repeated symbols denote the literal
ordered exponents, not powers of a map after discarding its original
intermediate positions.

\textbf{Theorem 2.} All pairs (6) satisfy
\SourceMath{n(v)\xrightarrow{l_\sigma}y(v)\xleftarrow{w_\sigma}m(v),
\qquad 0<m(v)<n(v),\qquad
\boxed{\nu_3(n(v))=b+t_e.}}{8} Both paths have exactly their displayed
valuations. Conversely the equality of these two original word
endpoints, on their positive sources, gives (5)-(6). The parameter
inverse is (n-n\_0)/(d\_sigma Q), with the same value recovered from m
or y using (6).

\textbf{Proof.} From (5), reduction modulo 3\^{}b first gives 3\^{}b
dividing n.~Write z=n/3\^{}b. Then Jz+h\_e=0 modulo 3\^{}a. Since
a\textgreater t\_e and J is a ternary unit, \(\nu_3(z)=t_e\), proving
the actual layer in (8). The integer \SourceMath{W=\frac{Jz+h_e}{3^a}}{}
is positive and has dyadic valuation one: J is divisible by 8, h\_e=2
modulo 4, and division by 3\^{}a preserves valuation one. Set
m=2\^{}aW-1. After j of the first a original returns the value is
\SourceMath{x_j=3^j2^{a-j}W-1\quad(0\le j\le a).}{9} For j\textless a
the valuation of x\_j+1 is at least two; therefore the next exponent is
exactly one. At j=a, the value is
\SourceMath{x_a=\frac{J}{3^b}n+h_e-1.}{} The identity 3h\_e-2=2\^{}e
gives the exact next exponent e and endpoint
\SourceMath{1+\frac{2^{2b-\sigma}n}{3^{b-1}}.}{} After j of the next b-1
steps the value is
\SourceMath{1+\frac{2^{2(b-j)-\sigma}n}{3^{b-1-j}}\quad(0\le j\le b-1).}{10}
At every nonterminal position of (10), its numerator correction is
divisible by 8; hence the next exponent is exactly two. The last value
is 4n+1 at sigma=0, and 2n+1 at sigma=1.

For sigma=0 the final exponent three gives (3n+1)/2, the left endpoint.
For sigma=1 the remaining exact path is
\SourceMath{2n+1\xrightarrow1 3n+2\xrightarrow1(9n+7)/2
\xrightarrow3(27n+23)/16.}{} Its valuations follow from n=27 modulo 32.
This is the left word's endpoint. All values are positive. Reversing the
original affine equality gives \(3^{a+b}W=Jn+h_e3^b\) and hence (5); the
original left valuations give the dyadic residue. This proves
completeness and both parameter inverses.

For height, retain the full affine identity
\SourceMath{m=(K/Q)n+h_e(2/3)^a-1.}{11} The coefficient K/Q is less than
one. Also (3) implies \SourceMath{h_e(2/3)^a<2^{\sigma-e}h_e(3/4)^b
\le \frac{2+2^{2-e}}{3}(3/4)^b<1.}{} Thus the constant in (11) is
negative and m\textless n.~Positivity follows from m=2\^{}aW-1 with
W\textgreater=2. QED.

No assertion that the different e-families are disjoint is made. Their
overlaps are calculated by the original congruences or by the full list
at each source. All eligible labels remain in that list, even when only
one is selected.

\section{4. Removing the apparent infinite search}

For a fixed admitted source n, the largest leading run in the same
template is \SourceMath{A_{\max}=\nu_3\left(J(n/3^b)+h_e\right).}{12}
Every original a' in the exact interval a\textless=a'\textless=A\_max
gives a smaller positive source
\SourceMath{m_{a'}=2^{a'}\frac{J(n/3^b)+h_e}{3^{a'}}-1.}{13} For
adjacent values,
\SourceMath{m_{a'+1}+1=\frac23(m_{a'}+1),\qquad T(m_{a'+1})=m_{a'}.}{}
The shared-future endpoint is unchanged. The rise endpoint is the same
integer \(J(n/3^b)+h_e-1\); the later right-hand path is unchanged as
well. Thus (12) selects the smallest source within this entire admitted
leading-run template. The exact signed chain correction is obtained by
adjoining the original incoming ones, as in the preceding bilateral
note.

\textbf{Theorem 3.} The complete set of admitted coefficient-contracting
templates at any positive odd source n is determined by finitely many
explicit tests. For n\textgreater1 put
\SourceMath{k(n)=\lfloor\log_2(n-1)\rfloor-1,
\qquad B=\nu_3(n).}{} Only even e with
\SourceMath{\boxed{2\le e\le k(n)}}{14} need be tested. For each e, the
only possible tail layer is \SourceMath{\boxed{b=B-\nu_3(e-1)\ge1.}}{15}
There are at most two support selectors to test at that (b,e). The tests
are the original dyadic residue, a\textless=k(n), and (5). Thus the
number of template labels examined is at most
\(2\lfloor k(n)/2\rfloor\), regardless of how large n is.

\textbf{Proof.} Every admitted smaller source has m=2\^{}\{a'\}W-1 with
W\textgreater=2 and m\textless n. Both sources are odd, so
m+1\textless=n-1. Hence \(2^{a'+1}\le n-1\) and \(a'\le k(n)\). Since
a'\textgreater=a\textgreater=e, (14) follows. Equation (8) and Lemma 1
give (15), uniquely. Every admitted template is therefore in the finite
list. Conversely, a successful test of its defining congruences is in
Theorem 2's complete source domain; (12) then gives the maximal run and
smallest template source. These implications prove the exact finite
search, rather than merely a termination heuristic. QED.

This proof handles all even exponents; it does not supply a bound on
arbitrary Collatz stopping times. It eliminates an unnecessary infinite
parameter search within the stated construction.

\section{5. A whole new root progression}

Take b=2,e=8,sigma=0. Then h=86, t\_e=0, a=16, and the original
common-future pair is \SourceMath{\boxed{\begin{aligned}
20\,241\,207+1\,549\,681\,956v
&\xrightarrow{1}30\,361\,811+2\,324\,522\,934v\\
&\xleftarrow{1^{16},8,2,3}
14\,024\,703+1\,073\,741\,824v,
\qquad v\ge0.
\end{aligned}}}{16} The exact height difference is
\SourceMath{6\,216\,504+475\,940\,132v>0.}{} The left sources have
actual ternary valuation two and no incoming original edge: 3x+1=2\^{}an
with 3 dividing n is impossible modulo 3. Both original paths are
needed. At v=0 the right path rises to 9,211,998,293 and then reaches
the displayed common endpoint. The comparison's signed odd-return clock
is 1-19=-18.

The complete previous reduction leaves every source of (16) as a root.
This is a finite exact residue certificate, not an extrapolation from
its anchor. The old finite mask has period 5,668,704. Since the new
progression has fixed ternary valuation, the four possible old e=2,6
layer congruences have fixed moduli. Their least common multiple with
the finite-mask period makes the membership of n\_0+n\_step v periodic
in v with period 8. The checker tests all eight parameter classes
against the full previous mask and all four original congruences. All
eight are roots. Periodicity proves the assertion for every
v\textgreater=0.

Consequently (16) is additional source coverage beyond both the
preceding finite catalogue and its infinite e=2,6 families. At its
anchor the new comparison ends at the still-retained root 14,024,703. No
arrival-at-1 conclusion is inferred from that class equality.

The other supported source orders are retained as well. For
b=1,e=4,sigma=1 the actual layer is two, not one, and the entire pair is
\SourceMath{45\,243+69\,984v\longrightarrow76\,349+118\,098v
\longleftarrow42\,367+65\,536v.}{} This illustrates the importance of
(2); no claim that this smaller family is new relative to the old finite
catalogue is made.

\section{6. Budget-preserving and unrestricted global retractions}

The primary new retraction first applies every old rule in its original
priority. Only at a retained root does it examine the complete finite
list in Theorem 3. For each candidate it constructs BOTH original paths,
their integral chain (1), and their maximum original vertex. It admits
the candidate only when that maximum is at most the predecessor's proved
budget
\SourceMath{R_0(n)=\max\{130(n+1),\lfloor64n(4/3)^{\lfloor\log_3n\rfloor}\rfloor+21\}.}{17}
Among the admitted candidates it selects the smallest target, with the
explicit word-length and label tie-breakers used in
\nolinkurl{all_even.py}. Rejected budget faces are not erased from the
full arithmetic candidate list. A source without an accepted move
remains a root.

All selected targets are strictly smaller odd integers. Therefore the
procedure terminates on every original input. For a selected move with
chain B\_n, target m, and signed clock j, define
\SourceMath{H(V_n)=B_n+q^jH(V_m),\qquad Q(V_n)=q^jQ(V_m),
\qquad F=I-Hd_\epsilon.}{} At roots use H=0 and Q=I. The exact
identities are \SourceMath{d_\epsilon H=I-Q,\quad Q^2=Q,\quad HQ=0,
\quad d_\epsilon F=Qd_\epsilon,\quad F^2=F,\quad FH=0.}{18} The first
telescopes (1). The next two follow at roots. Substitution in F proves
the remaining identities. Thus these are an actual finite-column
chain-homotopy retraction, not merely an equality of module ranks.

\textbf{Theorem 4.} The budget-preserving extension still uses only
original vertices at most R\_0(n) in its H column. For n\textgreater=27
the inherited bound \SourceMath{R_0(n)\le64n^{\log_3 4}+21}{} therefore
remains valid.

\textbf{Proof.} Every chosen reduction source x is at most the original
n.~Old moves have the original bound. The new acceptance test verifies
that bound for both complete original paths at x. The integer function
R\_0 is nondecreasing, so their union stays below R\_0(n). This also
handles cancellation: support is checked on both paths before equal edge
coefficients are collected. QED.

The whole new progression (16) passes this budget. For every n there,
n\textgreater=20,241,207\textgreater3\^{}15. Its largest right-hand
value is \SourceMath{(4096/9)n+85.}{} The inequality
\SourceMath{(4096/9)n+85\le64(4/3)^{15}n+21}{} holds at n=20,241,207 by
integer multiplication and has strictly positive margin slope. It holds
throughout the progression. The floor in (17) does not change this
integer endpoint conclusion. The short left path has smaller values.
This proves the budget acceptance on the entire family, not only at the
eight residue representatives used to test novelty.

For completeness, an unrestricted mode admits every successful template
from Theorem 3. For positive odd n\textgreater1 its larger support
budget is separately proved:
\SourceMath{\boxed{R_*(n)=\max\{130(n+1),\lfloor(n-1)3^{k(n)}/2^{k(n)}\rfloor-1\}.}}{19}
At the fixed root n=1, set R\_*(1)=1; the homotopy column H(V\_1) is
zero. During the initial rising run, m+1\textless=n-1 and
a'\textless=k(n), so \(3^{a'}W-1=(3/2)^{a'}(m+1)-1\) is bounded by the
second term. The later falling steps decrease; the bilateral final
three-edge segment stays below 5(n+1), covered by the first term. The
finite old paths also obey that first term. Hence (19) holds along the
entire decreasing-source recursion. In real exponents its second term is
at most \((2/3)(n-1)^{\log_2 3}-1\). This is not the smaller budget
(17), and it is not assigned that smaller exponent without the explicit
acceptance test.

Both versions have the integer uncollected-term bound, for positive odd
n\textgreater1,
\SourceMath{K_*(n)=\frac{n-1}{2}\max\{24,k(n)+\lfloor\log_3n\rfloor+6\}.}{20}
Set K\_*(1)=0 at the fixed root. There are at most (n-1)/2 strict source
reductions. A new pair uses at most k(n)+floor(log\_3 n)+6 edges in
total, and the finite older pairs use at most 24. Every signed
q-exponent is bounded in absolute value by that total. Therefore
collected constant coefficients have total absolute sum at most K\_*(n),
and first-order coefficients at most K\_*(n)\^{}2.

The old route is unchanged until its retained root. Its new tail gives
the comparison K=H\_new-H\_old, with
\SourceMath{d_\epsilon K=Q_{old}-Q_{new},\quad F_{old}K=K,
\quad F_{new}F_{old}=F_{old}F_{new}=F_{new}.}{21} The correction lies in
the old projected edge module: its boundary endpoints are old roots, on
which H\_old is zero. This proves the displayed F\_old identity;
substitution gives both projection compositions. The analogous Q
compositions follow from the fixed old route and the new root endpoints,
including all signed q-factors.

All arithmetic rule choices and budget tests are independent of epsilon.
Hence constant reduction commutes with these maps and transports the
integral first-jet comparison kernel through the actual homotopy
equivalence. The fixed loop stays present. Original periods are not
divided out, and a still-open component is not called absent because a
particular comparison sends an amplitude to supported zero.

\section{7. Verification and the still-global task}

The exact checker covers both selectors, b=1 through 16, and every even
e from 2 through 32. It verifies both original affine coefficient
columns at every step, then checks actual path instances, maximal
leading-run recovery, ternary layers, first-jet boundaries, and support
bounds. These are regression fixtures for the written all-parameter
proofs. The finite-search implementation is compared with a larger
direct template search on all odd inputs below 4,000. Retraction tests
verify (18) and (21) on original integral vectors, including newly
reduced roots. The separate affine auditor imports neither the new
module nor its predecessor.

The next arithmetic task is the exact residual root set after these
additional families and their budget faces. The finite algorithm for
template membership does not prove that one of its templates applies to
every retained source. The unchanged first-crossing horizon remains
11,610 odd steps in this increment; no extra crossing-window scan or
numerical stopping-time record is claimed.

A complete original singleton primitive would still require a route to
the fixed component, or another proved arithmetic exhaustion of the
surviving comparison kernel. New height relations improve its presented
source; they do not erase it. The cumulative archive keeps the earlier
source files and execution receipts separately from the checks executed
for this note.

\section{Sources and exact provenance}

Research direction and Split-Zero framework: the owner's workbench. This
is a continuation developed in the present tool-assisted session, not an
external independent mathematical review or a priority claim.

The mathematical predecessor is
\nolinkurl{bilateral_root_bounds_20260916/note.md}, with the earlier
original source and first-jet conventions preserved. The executable
predecessor is the unchanged \nolinkurl{bilateral_control.py}, whose
SHA-256 is pinned in \nolinkurl{SOURCE_INTAKE.json}. Current GitHub
integration and PR state were read on 16 September 2026 at main
\nolinkurl{9453b5306b3a897389dece309f9c119dffbcc489}; PR \#2 is now
merged. No remote write is inferred from generating this local addition.

For historical terminology, the coefficient-stopping-time question is
distinct from the present template-exhaustion problem. The original
expansion of an odd-return exponent a to one odd shortened step followed
by a-1 divisions is the clock comparison used in earlier notes.
Literature attribution remains in those notes and in Rozier--Terracol,
\emph{Paradoxical behavior in Collatz sequences},
\nolinkurl{https://arxiv.org/abs/2502.00948.} That paper supplies no
premise of Theorems 2--4.
